\documentclass[11pt,twoside]{article}

\usepackage[margin=1in]{geometry}
\usepackage{amsmath,amssymb,amsthm,mathtools}
\usepackage[inline,shortlabels]{enumitem}
\usepackage{booktabs}
\usepackage{xcolor}
\usepackage{microtype}

% Revision date: maintained by scripts/build_lecture_notes.py. Keep this command.
\newcommand{\lecturerevised}{2026-09-20 11:32 UTC-06:00}
\usepackage{eso-pic}
\AddToShipoutPictureBG{%
  \AtPageLowerLeft{%
    \put(\LenToUnit{0.5\paperwidth},\LenToUnit{0.5in}){%
      \makebox(0,0){\normalfont\footnotesize\color{black!60}Last revised: \lecturerevised}%
    }%
  }%
}

\usepackage[square,authoryear]{natbib}
\usepackage[colorlinks=true,linkcolor=blue!60!black,
  urlcolor=blue!60!black,citecolor=blue!60!black]{hyperref}
\usepackage[nameinlink,capitalize]{cleveref}
\hypersetup{pdftitle={Lecture 6: Data collection, dependence, and useful evidence},
  pdfauthor={Csaba Szepesvari}}
\crefname{exercise}{Exercise}{Exercises}
\Crefname{exercise}{Exercise}{Exercises}

\usepackage{enotez}
\DeclareInstance{enotez-list}{lecture}{list}{format=\normalsize,number={#1.}}
\setenotez{list-name=Endnotes,list-style=lecture,backref=true}
\crefname{endnote}{Endnote}{Endnotes}
\Crefname{endnote}{Endnote}{Endnotes}

\definecolor{courseblue}{HTML}{17365D}
\newcommand{\E}{\mathbb E}
\newcommand{\I}{\mathbb I}
\renewcommand{\P}{\mathbb P}
\newcommand{\R}{\mathbb R}
\newcommand{\KL}{D}
\newcommand{\TV}{\mathrm{TV}}
\makeatletter
\@for\@letter:=A,B,C,D,E,F,G,H,I,J,K,L,M,N,O,P,Q,R,S,T,U,V,W,X,Y,Z\do{%
  \expandafter\edef\csname c\@letter\endcsname{\noexpand\mathcal{\@letter}}}
\makeatother
\newtheorem{theorem}{Theorem}[section]
\newtheorem{proposition}[theorem]{Proposition}
\newtheorem{corollary}[theorem]{Corollary}
\theoremstyle{definition}
\newtheorem{definition}[theorem]{Definition}
\newtheorem{example}[theorem]{Example}
\theoremstyle{remark}
\newtheorem{remark}[theorem]{Remark}

\renewcommand{\thepage}{6--\arabic{page}}
\renewcommand{\thesection}{6.\arabic{section}}
\renewcommand{\theequation}{6.\arabic{equation}}
\pagestyle{myheadings}
\markboth{Lecture 6: Data collection, dependence, and useful evidence}
{Lecture 6: Data collection, dependence, and useful evidence}
\usepackage{../exercises/course-exercises}

\begin{document}

\begin{center}
\fcolorbox{courseblue}{white}{%
\begin{minipage}{0.92\textwidth}
\vspace{1mm}
\textbf{CMPUT 654: Mathematical Foundations of Modern AI Systems}
\hfill Fall 2026

\vspace{3mm}
\begin{center}
{\Large\bfseries Lecture 6: Data collection, dependence, and useful evidence}

\vspace{1mm}
September 17, 2026
\end{center}

\vspace{2mm}
\textit{Lecturer: Csaba Szepesv\'ari}
\hfill
\textit{Note prepared by: Csaba Szepesv\'ari}
\vspace{1mm}
\end{minipage}}
\end{center}

\noindent Larkin Wisdom prepared a scribe of the September 17 lecture.
Csaba Szepesv\'ari prepared and expanded the present notes.

\medskip
\noindent\textbf{Reading guide.}
These notes ask how collected data become evidence about deployment. The main
text follows the classroom discussion: corpus construction, the role of
randomness, concentration in the iid model, exact copies, resampling, and
beta-mixing. The public pretraining references, the longer resampling
calculation, and the proof behind the mixing theorem were added after class.
The exercises develop calculations that were posed but not completed in class.
Lectures 1--5 provide the background on sequence prediction, loss, and
state-space models.

\section{A dataset is produced by a procedure}
\label{sec:pipeline}

Data is a major part of an ML system. The iid model deliberately suppresses
most details of how data is obtained. Before using that model, we should ask
where the data came from and what happened to it.

Language-model pretraining commonly starts from web pages, books, code,
scientific articles, conversations, and other text collections. A simplified
pipeline is
\[
 \begin{gathered}
 \text{sources}
 \longrightarrow \text{extraction and parsing}
 \longrightarrow \text{filtering and deduplication}\\
 \longrightarrow \text{source mixture}
 \longrightarrow \text{tokenization}
 \longrightarrow \text{training sequences}.
 \end{gathered}
\]
This pipeline raises a recurring question: does deduplication matter, when does
it matter, and by how much? \Cref{sec:copies} answers it under a simple variance
model.
Extraction may remove HTML and recover the main text. Filters may use language,
quality, safety, or document-type criteria. Exact and approximate
deduplication remove repeated material. A mixture rule determines how much of
each retained source is used. Tokenization turns text into token IDs.
Finally, the tokens must be placed into sequences of the length used for
training. Public descriptions of open-weight models document many, though not
always all, of these choices; the bibliographical notes give examples.

A common sequence-construction method concatenates documents, places a special
boundary token between them, and cuts the resulting stream into fixed-length
sequences. Other methods pack shorter documents without allowing attention
across document boundaries or truncate long documents. These choices determine
which prediction contexts training presents to the model. They can also make
training contexts unlike deployment contexts: a training sequence may usually
be full, whereas a deployed model must also predict after a short prompt.
In a \emph{causally masked transformer}, the prediction at a position cannot
use later tokens. A full training sequence therefore contains prediction
problems with prefix lengths ranging from short to long. This is a
transformer-specific mitigation, not a general answer to mismatch.

The guiding principle is that training data must resemble deployment data in
the respects needed for the desired behavior. Exact equality of distributions
is not required, and post-training can alter behavior. It cannot generally
recover abilities for which pretraining and post-training supply no relevant
evidence. Pretraining only on Hungarian documents and deploying only in
English preserves some common structure, but not enough to expect good English
behavior.

Three objects should remain separate:
\begin{enumerate}[leftmargin=*]
\item the \emph{realized corpus}, which is a fixed collection of documents and
  sequences;
\item the \emph{collection procedure}, which found, selected, filtered,
  copied, and divided the sources; and
\item a \emph{probability model}, which describes repeated collection, an
  ordered source, evaluation cases, or future deployment.
\end{enumerate}
A list of files does not determine a probability model. The proposed unit---a
token, window, document, conversation, user, or source---also matters. One book,
one thousand independently written books, and one book stored one thousand
times can all yield one thousand training sequences but provide different
evidence.

\section{Why introduce randomness?}
\label{sec:randomness}

A probability law specifies which future cases count and how they are
weighted. A joint law additionally states how observations are related. It can
therefore support quantitative claims about whether empirical averages are
stable and whether they estimate deployment quantities.

Randomness can enter in three distinct places:
\begin{itemize}[leftmargin=*]
\item \emph{collection randomness}: units are sampled from a stated
  population or by a randomized experiment;
\item \emph{source randomness}: a stochastic process models future documents,
  interactions, or tokens; and
\item \emph{algorithmic randomness}: shuffling, minibatch selection,
  initialization, and decoding randomize computation.
\end{itemize}
Algorithmic randomness does not turn repeated source material into independent
population evidence. Randomness may also be an effective mathematical
description of a deterministic system, but establishing such a description is
a separate problem.\endnote{Deterministic systems can exhibit statistical
regularities. Statistical mechanics, dynamical systems, and pseudorandomness
give precise versions in particular settings. Their existence does not by
itself justify a stochastic model for a corpus.}\label{en:deterministic-randomness}

Observations that are not independent are called \emph{dependent}. An iid
model can be plausible when, for example, a machine independently selects
problems from a fixed collection or unrelated users independently initiate
conversations under stable conditions. These remain modeling claims about the
chosen units and procedure.

Independence is also not necessary for learning. For the threshold class
\(f_k(x)=\I\{x\geq k\}\), the two labeled examples
\((k-1,0)\) and \((k,1)\) identify an interior threshold exactly. This is a
deterministic teaching set: the teacher knows the target and chooses examples
that identify it. This differs from active learning, where the learner does not
know the target but may choose which labels to request, and from passive
learning, where the learner receives examples from a fixed sampling law. For
interior thresholds on a finite ordered set, teaching uses two target-dependent
examples, active learning permits binary search, and passive learning must wait
for the sampling law to reveal points near the threshold. A boundary threshold
can be taught with one example. Dependence alone does not explain these
differences; the information available when examples are chosen does.

Expert curation for language models can resemble teaching when examples are
selected to communicate a desired capability. The analogy concerns data
selection: a target-aware example can be more informative than a passive draw.
It does not imply that one or a few examples suffice for a model to learn or
transfer a proof. The model class, optimization procedure, and relation between
training and deployment tasks still matter.

A dependent process can also provide quantitative coverage. Suppose the
deployment law on \([N]\) is \(P\), and, for every interval \(I\subseteq[N]\),
\begin{equation}
 \P(X_i\in I\mid X_1,\ldots,X_{i-1})\geq \alpha P(I)
 \quad\text{almost surely}.
 \label{eq:conditional-coverage}
\end{equation}
For the threshold learner from Lecture 2, let \(L\) denote the deployment
probability of its error interval. The event \(L>\epsilon\) requires the sample
to miss a fixed interval, determined by \(P\), the target, and \(\epsilon\),
whose \(P\)-mass exceeds \(\epsilon\). Each observation has conditional
probability at least \(\alpha\epsilon\) of hitting that interval. Iterated
conditioning gives
\begin{equation}
 \P(L>\epsilon)
 \leq (1-\alpha\epsilon)^n
 \leq \exp(-\alpha n\epsilon).
 \label{eq:conditional-coverage-bound}
\end{equation}
Thus independence is one sufficient explanation for accumulating evidence,
not the definition of learnability.

\section{The iid calculation}
\label{sec:iid}

Let \(X_1,\ldots,X_n\) be iid with law \(P\). For a fixed parameter
\(\theta\), define the \emph{population loss} and \emph{empirical loss} by
\begin{equation}
 L(\theta)=\E[\ell(\theta;X)],
 \qquad
 \widehat L_n(\theta)=\frac1n\sum_{i=1}^n\ell(\theta;X_i).
 \label{eq:risks}
\end{equation}
If \(0\leq\ell(\theta;X)\leq1\), Hoeffding's inequality gives
\begin{equation}
 \P\!\left(\left|\widehat L_n(\theta)-L(\theta)\right|>\epsilon\right)
 \leq 2e^{-2n\epsilon^2}.
 \label{eq:hoeffding}
\end{equation}
This is a \emph{concentration} statement: the empirical average is unlikely to
be far from its expectation. The systematic study of such behavior is called
measure concentration.

The variance explains the first-order scale. Writing
\(\sigma_\theta^2=\operatorname{Var}(\ell(\theta;X))\), independence gives
\begin{equation}
 \operatorname{Var}(\widehat L_n(\theta))
 =\frac{\sigma_\theta^2}{n}.
 \label{eq:iid-variance}
\end{equation}
When \(0<\sigma_\theta^2<\infty\), the central limit theorem further gives
\[
 \sqrt n\bigl(\widehat L_n(\theta)-L(\theta)\bigr)
 \ \Longrightarrow\ \mathcal N(0,\sigma_\theta^2).
\]
Equivalently, for large \(n\), the deviation
\(\widehat L_n(\theta)-L(\theta)\) is approximately
\(\mathcal N(0,\sigma_\theta^2/n)\).
These statements concern a fixed \(\theta\). Applying them to a parameter
chosen from the same data requires a uniform, stability, information, or other
generalization argument.

The copied process \(Y_i=U\), where
\(U\sim\operatorname{Bernoulli}(1/2)\), shows what independence prevents.
Every \(Y_i\) has the same marginal law, but
\(n^{-1}\sum_iY_i=U\) for every \(n\). The empirical average does not
concentrate around \(1/2\).

\section{Exact copies and effective sample size}
\label{sec:copies}

We can now answer the question raised at the start of class: does
deduplication matter, when does it matter, and by how much?

Let \(X_1,\ldots,X_m\) be independent source units with common loss variance
\(\sigma_\theta^2\). Suppose unit \(i\) occurs \(c_i\geq1\) times in the
stored corpus, and let \(N=\sum_{i=1}^m c_i\). The stored-data average is
\begin{equation}
 \widehat L_{\rm copy}(\theta)
 =\frac1N\sum_{i=1}^m c_i\ell(\theta;X_i).
 \label{eq:copy-average}
\end{equation}
Independence of the \(m\) source units gives
\begin{equation}
 \operatorname{Var}(\widehat L_{\rm copy}(\theta))
 =\sigma_\theta^2\frac{\sum_{i=1}^m c_i^2}{N^2}
 =\frac{\sigma_\theta^2}{n_{\rm eff}},
 \qquad
 n_{\rm eff}=\frac{N^2}{\sum_{i=1}^m c_i^2}.
 \label{eq:copy-neff}
\end{equation}
The \emph{effective sample size} is the independent sample size that would
give the same variance under this model. Cauchy--Schwarz gives
\begin{equation}
 1\leq n_{\rm eff}\leq m,
 \label{eq:neff-bounds}
\end{equation}
with \(n_{\rm eff}=m\) exactly when all multiplicities are equal. If one
source dominates, \(n_{\rm eff}\) can be close to one.

Deduplicating to one copy of every source gives variance
\(\sigma_\theta^2/m\). Hence
\[
 \frac{\operatorname{Var}(\widehat L_{\rm copy})}
      {\sigma_\theta^2/m}
 =m\frac{\sum_i c_i^2}{N^2}
 =1+\operatorname{CV}(c_1,\ldots,c_m)^2,
\]
where \(\operatorname{CV}\) is the empirical coefficient of variation using
denominator \(m\). Equal duplication wastes storage and computation but does
not reduce effective sample size below \(m\). Unequal duplication also reduces
effective sample size and changes the empirical weights. If multiplicity is
intended to represent deployment frequency, deduplication changes the target;
Lecture 7 studies that weighting question.

Near copies require the more general covariance calculation below. Their value
depends on the statistic: two variants can be nearly identical for one learned
feature and materially different for another.

\section{Covariance is the general calculation}
\label{sec:covariance}

Let \(G_t\) be a scalar loss, score, or projection of a gradient contribution
at position \(t\). No independence assumption is needed for
\begin{equation}
 \operatorname{Var}\!\left(\frac1T\sum_{t=1}^T G_t\right)
 =\frac1{T^2}\sum_{s,t=1}^T\operatorname{Cov}(G_s,G_t).
 \label{eq:covariance-sum}
\end{equation}
Define
\[
 S_T=\sum_{t=1}^T\operatorname{Var}(G_t),
 \qquad
 \Gamma_T=\sum_{s,t=1}^T\operatorname{Cov}(G_s,G_t).
\]
When \(S_T>0\) and \(\Gamma_T>0\), variance matching gives the
statistic-dependent effective sample size
\begin{equation}
 T_{\rm eff}(G)=T\frac{S_T}{\Gamma_T}.
 \label{eq:general-neff}
\end{equation}
Independence gives \(T_{\rm eff}=T\). Positive covariance usually reduces it;
negative covariance can increase it. There is no single effective sample size
for a corpus without specifying the quantity being averaged.

This distinction matters for sequences. A single long trajectory may contain
many effective observations of a shared transition parameter but only one
draw from its initial-state law. A book may repeatedly inform local syntax
while retaining the same notation convention throughout. Overlapping context
windows share most of their tokens; \cref{ex:l6-overlapping-windows} asks for
their exact dependence range in a simple model.

\section{What resampling changes}
\label{sec:resampling}

Fix a realized corpus \(x_1,\ldots,x_m\). Sampling indices independently with
replacement produces conditionally independent draws from its empirical law.
It does not create new independent evidence about the population that produced
the corpus. As the number of resampled draws grows with the corpus fixed, the
random frequencies approach the original empirical frequencies, so the same
empirical average is recovered.

Sampling all \(m\) indices without replacement is a random shuffle. Sampling
\(B<m\) without replacement is equivalently a random shuffle followed by
truncation. Random minibatches can change the optimization path, but repeated
minibatch training still targets the original empirical objective. Sparse
subsampling may separate nearby observations in a dependent sequence, at the
cost of discarding data and adding variability. The appendix makes these
claims precise; \cref{ex:l6-resampling} asks you to compare sampling with and
without replacement.

\section{Controlled dependence through beta-mixing}
\label{sec:mixing}

There are many ways to quantify dependence. For a one-sided process
\(X_1,X_2,\ldots\), let \(\mathcal L(Y)\) denote the distribution of a
random element \(Y\). Define the \emph{absolute-regularity}, or
\emph{beta-mixing}, coefficient
\begin{equation}
 \beta(a)=\sup_{t\geq1}
 \left\|
 \mathcal L(X_{1:t},X_{t+a:\infty})
 -\mathcal L(X_{1:t})\otimes\mathcal L(X_{t+a:\infty})
 \right\|_{\TV},
 \qquad a\geq1.
 \label{eq:beta}
\end{equation}
Here \(\|P-Q\|_{\TV}=\sup_A|P(A)-Q(A)|\). Small \(\beta(a)\) means that an
observed past and a future separated from it by a gap of \(a-1\) positions are
close to independent in a strong, distributional sense.

Select positions \(t_j=1+(j-1)a\), for \(j=1,\ldots,q\), where
\(q=1+\lfloor(n-1)/a\rfloor\). Assume the selected observations have common
marginal law \(P\). The independent-points comparison proved in
\cref{app:mixing-proof} is
\begin{equation}
 \left\|\mathcal L(X_{t_1},\ldots,X_{t_q})-P^{\otimes q}\right\|_{\TV}
 \leq(q-1)\beta(a).
 \label{eq:independent-points}
\end{equation}

\begin{theorem}[Finite-class consistency under beta-mixing]
\label{thm:mixing-finite-class}
Let \(\cF\) be finite, let the target \(f\in\cF\), and, for \(X\sim P\),
define \(L_P(g,f)=P(g(X)\neq f(X))\). If a learner returns a function
\(\widehat f\in\cF\) consistent with all labeled observations
\((X_i,f(X_i))\), \(i\in[n]\), then
\begin{equation}
 \P\bigl(L_P(\widehat f,f)>\epsilon\bigr)
 \leq |\cF|(1-\epsilon)^q+(q-1)\beta(a).
 \label{eq:mixing-bound}
\end{equation}
\end{theorem}

\begin{proof}
Let \(E\) be the event that some \(g\in\cF\) with
\(L_P(g,f)>\epsilon\) is consistent at every selected position. Full-sample
consistency implies the failure event is contained in \(E\). By
\cref{eq:independent-points} and the iid finite-class argument from Lecture 2,
\[
 \P(E)\leq P^{\otimes q}(E)+(q-1)\beta(a)
 \leq |\cF|(1-\epsilon)^q+(q-1)\beta(a).
\]
\end{proof}

Larger gaps reduce the dependence term but leave fewer selected observations.
For example, if \(\beta(a)\leq e^{-a}\), choosing
\(a=\lceil\log(2n/\delta)\rceil\) makes the dependence term at most
\(\delta/2\). The remaining term behaves as though the sample size were
roughly \(n/\log(n/\delta)\): this blocking argument loses only a logarithmic
factor compared with iid sampling. \Cref{ex:l6-polynomial-mixing} asks what the
same argument gives under polynomial decay.

Finite irreducible aperiodic Markov chains mix geometrically. On a compact
Polish state space, irreducibility and the strong Feller property make the
whole state space a small set; with an appropriate aperiodicity condition,
this gives a Doeblin route to geometric ergodicity and hence exponential
beta-mixing. On a noncompact state space, the same argument only makes compact
sets small, so a return or drift condition is generally needed. Ergodic chains
can mix polynomially or more slowly. The point here is the distinction between
the clean compact case and the additional control needed on noncompact spaces,
not a catalogue of minimal conditions.

Mixing is highly relevant when learning from one long reinforcement-learning
trajectory or another continuing process. It is not automatically a realistic
global model of a web corpus. A book can introduce notation at the beginning
and use it to the end. An author, topic, task, or other slowly evolving state
can create arbitrarily long unconditional dependence even if local residuals
forget quickly after conditioning on that state. A state-space account with
slow and fast components may then be more informative. Mixing must always be
checked for the process and statistic to which a theorem is applied.

\section{Dependence is also predictive information}
\label{sec:signal}

The prefix--target pairs in one document are dependent and strongly overlap.
That dependence is not merely a defect: earlier tokens contain information
about later ones. A sequence model reuses parameters across positions so that
many different contexts can constrain shared predictive structure. The same
trajectory can therefore provide useful evidence without becoming a collection
of iid examples.

Parameter sharing alone does not prove that the shared regularity is correct,
that training covers deployment contexts, or that optimization recovers it.
The relevant question is whether the observed dependent data constrain the
particular shared object needed at deployment.

Evaluation uses the same logic. Define the intended deployment unit and source;
keep exact copies, near copies, and overlapping chunks from crossing the
training--test boundary; and remember that repeated selection on a test set
makes it part of development. A common law \(P\) for iid training and test
examples is the strongest simple stability model, not a default fact about
data.

\section{Take-home messages}
\label{sec:take-home}

\begin{itemize}[leftmargin=*]
\item A corpus is the output of a collection and processing pipeline. A
  probability model is an additional analytical choice.
\item Independence gives useful quantitative concentration, but deterministic
  teaching and dependent coverage show that it is not necessary for learning.
\item Exact duplication changes empirical weights and can reduce effective
  sample size. The variance formula states precisely when deduplication helps.
\item The covariance of the statistic being averaged determines how much
  evidence dependent observations provide.
\item Resampling or shuffling a fixed corpus does not create new population
  evidence.
\item Mixing and blocking recover weakened iid guarantees under controlled
  dependence. Persistent state can make unconditional mixing inappropriate.
\item Dependence within a sequence is predictive signal. Evaluation is useful
  only for the deployment claim represented by its construction.
\end{itemize}

\clearpage
\printendnotes

\section*{Bibliographical notes}

Public pretraining accounts illustrate which pipeline details are documented.
Dolma describes source-specific extraction, filtering, deduplication, and
mixture construction \citep{soldaini2024}. FineWeb reports large-scale web
filtering and deduplication studies \citep{penedo2024}. The Llama 3 report
describes its data mixture and filtering at a higher level
\citep{dubey2024}. Sequence packing and document-boundary choices are analyzed
explicitly by \citet{ding2024}.

Hoeffding's inequality is from \citet{hoeffding1963}. The independent-block
method for mixing empirical processes is developed by \citet{yu1994}.
The exercise contrasting the central limit theorem with exponential tail
bounds is adapted from the distinction developed in Exercises~5.10 and~5.13 of
\citet{lattimore2020}.
\citet{hairer2002} proves that every compact set is small for an irreducible
strong Feller chain. If the state space is compact, this applies to the whole
space. \citet{bradley2005} surveys strong mixing conditions and records that
geometric ergodicity of a stationary Markov chain implies exponential
beta-mixing. These results give the sufficient route stated in the text; they
are not presented as minimal conditions.

\section*{Glossary}

\begin{description}[style=nextline,leftmargin=1.5em,labelindent=0pt]
\item[Absolute regularity or beta-mixing]
The total-variation distance in \cref{eq:beta} between the joint law of a past
and a separated future and the product of their marginal laws.
\item[Collection procedure]
The operations that find, select, filter, copy, and divide sources into the
realized data.
\item[Concentration]
The property that a random quantity, such as an empirical average, is unlikely
to be far from a specified central value.
\item[Effective sample size]
The independent sample size whose variance matches that of a specified
dependent or unequally weighted average.
\item[Empirical loss]
The average loss on the realized training observations, as in
\cref{eq:risks}.
\item[Population loss]
Expected loss under a stated law representing the cases of interest.
\item[Realized corpus]
The fixed collection of documents, tokens, or sequences actually available.
\item[Resampling with replacement]
Independent draws from the empirical distribution of a fixed dataset.
\end{description}

\section{Exercises}
\label{sec:exercises}

\exerciseinstructions
\setcounter{courseexercise}{0}
\begin{exercisechunk}
\item \textbf{Teaching, active learning, and passive learning.}
\label[exercise]{ex:l6-threshold-information}
\exerciseprofile{2}{2}{2}{\exproof\quad\excalculation\quad\exinterpretation}{%
Distinguish three information structures by deriving their requirements for
identifying a threshold.}
Let \(N\geq6\), let \(\mathcal K=\{2,\ldots,N-1\}\), and, for
\(k\in\mathcal K\), define \(f_k(x)=\I\{x\geq k\}\) on \([N]\).
The goal is to identify \(k\) exactly from labeled examples.
\begin{enumerate}[(a),beginpenalty=10000]
\item A teacher knows \(k\) and may choose labeled examples. Prove that
\(\{(k-1,0),(k,1)\}\) identifies \(k\) uniquely within \(\mathcal K\).
Explain why no single labeled example identifies every \(k\in\mathcal K\).
\item An active learner does not know \(k\), but may choose each query after
observing the previous labels. Give a binary-search strategy that identifies
every \(k\in\mathcal K\) using at most
\(\lceil\log_2(N-2)\rceil\) queries. Prove that no active strategy can have a
smaller worst-case number of binary-label queries.
\item A passive learner receives \(X_1,\ldots,X_n\) independently and uniformly
from \([N]\), together with their labels. Fix \(k\in\{3,\ldots,N-2\}\).
Prove that the observed labels identify \(k\) uniquely if and only if both
\(k-1\) and \(k\) occur in the sample, and that the probability of this event is
\[
 1-2\left(1-\frac1N\right)^n
   +\left(1-\frac2N\right)^n.
\]
Deduce that \(n\geq N\log(2/\delta)\) suffices for this probability to be at
least \(1-\delta\).
\item Explain why the two-example teaching set is not evidence that dependence
alone makes learning faster. State precisely what target information the
teacher has that the active and passive learners do not.
\end{enumerate}
\begin{hints}
For (b), maintain the set of thresholds consistent with all answers. For the
lower bound, a depth-\(q\) binary decision tree has at most \(2^q\) leaves. For
(c), use inclusion--exclusion and then a union bound.
\end{hints}
\end{exercisechunk}

\begin{exercisechunk}
\item \textbf{Hoeffding bounds and the normal approximation.}
\label[exercise]{ex:l6-hoeffding-clt}
\exerciseprofile{2}{2}{2}{\excalculation\quad\exinterpretation}{%
Distinguish the central-limit scale from fixed-size deviations and compare the
conclusions supplied by the central limit theorem and Hoeffding's inequality.}
Let \(X_1,\ldots,X_n\) be iid \(\operatorname{Bernoulli}(p)\) variables with
\(p\in(0,1)\), and let \(\widehat p_n=n^{-1}\sum_{i=1}^nX_i\).
\begin{enumerate}[(a),beginpenalty=10000]
\item Derive \(\E[\widehat p_n]\) and
\(\operatorname{Var}(\widehat p_n)\), and state the central limit theorem for
\(\sqrt n(\widehat p_n-p)\).
\item Apply Hoeffding's inequality to show that, for every \(\epsilon>0\),
\[
 \P(\widehat p_n-p\geq\epsilon)\leq e^{-2n\epsilon^2}.
\]
\item Fix \(t>0\) and set
\(\epsilon_n=t\sqrt{p(1-p)/n}\). Determine the limit of
\(\P(\widehat p_n-p\geq\epsilon_n)\), and show that the Hoeffding bound in
part~(b) becomes \(e^{-2t^2p(1-p)}\). Explain what the two conclusions say at
the \(n^{-1/2}\) scale.
\item Now fix \(\epsilon\in(0,1-p)\), and let
\(Z_n\sim\mathcal N(p,p(1-p)/n)\). You may use
\[
 \lim_{x\to\infty}\frac{1}{x^2}\log(1-\Phi(x))=-\frac12.
\]
Show that
\[
 \lim_{n\to\infty}\frac1n
 \log\P(Z_n-p\geq\epsilon)
 =-\frac{\epsilon^2}{2p(1-p)}.
\]
Explain why the central limit theorem alone does not justify using this
Gaussian rate for \(\P(\widehat p_n-p\geq\epsilon)\), and contrast this with
the finite-sample conclusion in part~(b).
\end{enumerate}
\end{exercisechunk}

\begin{exercisechunk}
\item \textbf{Exact duplication and effective sample size.}
\label[exercise]{ex:l6-duplication}
\exerciseprofile{2}{2}{3}{\exproof\quad\excalculation\quad\exinterpretation}{%
Derive how unequal duplication changes variance and effective sample size, and
state what this calculation does not establish about learning performance.}
Let \(X_1,\ldots,X_m\) be independent, and suppose
\(Z_i=\ell(\theta;X_i)\) has variance \(\sigma_\theta^2>0\) for every
\(i\). Unit \(i\) is stored \(c_i\geq1\) times. Set
\(N=\sum_{i=1}^m c_i\) and
\(\widehat L=N^{-1}\sum_{i=1}^m c_iZ_i\).
\begin{enumerate}[(a),beginpenalty=10000]
\item Derive \(\operatorname{Var}(\widehat L)\) and the effective sample size
\(n_{\rm eff}\) in \cref{eq:copy-neff}.
\item Prove \(1\leq n_{\rm eff}\leq m\). State exactly when the upper bound
is attained.
\item Calculate \(n_{\rm eff}\) when (i) \(c_i=c\) for every \(i\), and
(ii) \(c_1=r\) and \(c_i=1\) for \(i\geq2\). Find the limit in (ii) as
\(r\to\infty\) with \(m\) fixed.
\item Let \(\bar c=N/m\) and
\(\operatorname{CV}^2=m^{-1}\sum_i(c_i-\bar c)^2/\bar c^2\). Prove
\[
 \frac{\operatorname{Var}(\widehat L)}{\sigma_\theta^2/m}
 =1+\operatorname{CV}^2.
\]
Explain what this ratio does and does not establish about deduplication.
\end{enumerate}
\begin{hints}
For (b), apply Cauchy--Schwarz to \(\sum_i c_i\). For the lower bound,
compare \((\sum_i c_i)^2\) and \(\sum_i c_i^2\).
\end{hints}
\end{exercisechunk}

\exercisegroup{Sampling from a finite corpus}
\begin{exercisechunk}
\item \textbf{Sampling with and without replacement.}
\label[exercise]{ex:l6-resampling}
\exerciseprofile{2}{2}{3}{\exproof\quad\excalculation\quad\exinterpretation}{%
Derive the laws induced by sampling indices with and without replacement and
identify what these procedures do not change about the source data.}
Fix a corpus \(x_1,\ldots,x_m\in\R\) and let
\(\bar x=m^{-1}\sum_i x_i\).
\begin{enumerate}[(a),beginpenalty=10000]
\item Draw \(I_1,\ldots,I_B\) independently and uniformly from \([m]\), and
define \(C_i=\sum_{b=1}^B\I\{I_b=i\}\). Give the joint law of
\((C_1,\ldots,C_m)\), and prove that
\(B^{-1}\sum_iC_ix_i\) is conditionally unbiased for \(\bar x\).
\item Take \(B=m\). Find the marginal law of \(C_i\) and its limit as
\(m\to\infty\). Find the limiting probability that item \(i\) is omitted.
Why does this calculation not show that the resampled items are independent
population observations?
\item Suppose instead that indices are sampled without replacement. Prove that
sampling all \(m\) indices gives a uniformly random permutation. Show that the
first \(B<m\) indices are a uniformly random ordered sample without replacement,
or equivalently a uniformly random \(B\)-element subset followed by a uniformly
random ordering.
\end{enumerate}
\end{exercisechunk}

\begin{exercisechunk}
\item \textbf{Dependence retained by resampling.}
\label[exercise]{ex:l6-resampling-dependence}
\exerciseprofile{2}{1}{1}{\exproof\quad\excalculation\quad\exinterpretation}{%
Use total covariance to distinguish conditional independence of resampled
draws from unconditional dependence inherited from the corpus.}
Let \(X_1,\ldots,X_m\) be square-integrable and possibly dependent. Independently of
\(X_{1:m}\), draw \(I_1,I_2\) independently and uniformly from \([m]\), and
set \(Y_b=X_{I_b}\). Prove
\[
 \operatorname{Cov}(Y_1,Y_2)=\operatorname{Var}\!\left(
   \frac1m\sum_{i=1}^mX_i\right).
\]
Interpret the result.
\begin{hints}
Condition on \(X_{1:m}\) and use the law of total covariance.
\end{hints}
\end{exercisechunk}

\begin{exercisechunk}
\item \textbf{Dependence created by overlapping windows.}
\label[exercise]{ex:l6-overlapping-windows}
\exerciseprofile{3}{2}{2}{\exproof\quad\excalculation\quad\exinterpretation}{%
Compute the beta-mixing coefficient induced by overlapping iid windows and
distinguish token independence from window independence.}
Let \(\mathcal Z\) be a finite or countable alphabet, let
\(p:\mathcal Z\to[0,1]\) satisfy \(\sum_{z\in\mathcal Z}p(z)=1\), and let
\((Z_t)_{t\geq1}\) be iid \(\mathcal Z\)-valued random variables with law
\(p\). Let \(L\geq2\) be an integer, and define the \(\mathcal Z^L\)-valued
windows \(W_t=(Z_t,\ldots,Z_{t+L-1})\). For an integer \(a\geq1\), define the
beta-mixing coefficient
\[
 \beta_W(a)=\sup_{t\geq1}
 \left\|\mathcal L(W_{1:t},W_{t+a:\infty})
 -\mathcal L(W_{1:t})\otimes\mathcal L(W_{t+a:\infty})\right\|_{\TV},
\]
\begin{enumerate}[(a),beginpenalty=10000]
\item Show that the window process is \((L-1)\)-dependent: collections of
windows separated by at least \(L\) starting positions are independent.
Deduce that \(\beta_W(a)=0\) for \(a\geq L\).
\item For \(1\leq a<L\), show that
\[
 \beta_W(a)=1-\left(\sum_{z\in\mathcal Z}p(z)^2\right)^{L-a}.
\]
Explain from this formula why the coefficient is close to one when the
collision probability \(\sum_{z\in\mathcal Z}p(z)^2\) is small or when the number \(L-a\)
of shared positions is large, even though the underlying token process is
independent.
\end{enumerate}
\begin{hints}
For (b), identify the underlying variables that occur in both window
collections.
\end{hints}
\end{exercisechunk}

\begin{exercisechunk}
\item \textbf{The cost of polynomial beta-mixing.}
\label[exercise]{ex:l6-polynomial-mixing}
\exerciseprofile{2}{2}{2}{\exproof\quad\excalculation\quad\exinterpretation}{%
Convert a polynomial mixing rate into a sample-size bound and compare its
dependence on accuracy with iid and geometric-mixing bounds.}
Assume the setting of \cref{thm:mixing-finite-class} and
\(\beta(a)\leq Ca^{-r}\) for constants \(C,r>0\). Let
\(\epsilon,\delta\in(0,1)\) and \(A=\log(2|\cF|/\delta)\).
\begin{enumerate}[(a),beginpenalty=10000]
\item Starting from \cref{eq:mixing-bound}, give sufficient conditions on
the selected-sample size \(q\) and gap \(a\) that make each term at most
\(\delta/2\).
\item Use \(n\geq1+(q-1)a\) and a suitable choice of \(q,a\) to show, up to
universal constants and integer rounding, that it suffices to take
\[
 n\gtrsim \frac{A}{\epsilon}
 +\left(\frac{A}{\epsilon}\right)^{1+1/r}
  \left(\frac{C}{\delta}\right)^{1/r}.
\]
\item Compare the dependence on \(\epsilon\) with iid sampling and with
geometric beta-mixing. Is the degradation exponential in \(1/\epsilon\)?
\end{enumerate}
\begin{hints}
Use \((1-\epsilon)^q\leq e^{-\epsilon q}\). First choose
\(q\asymp A/\epsilon\), then choose \(a\) so that
\(qCa^{-r}\lesssim\delta\), and finally take \(n\) of order \(qa\).
\end{hints}
\end{exercisechunk}


\appendix

\section{Resampling a fixed corpus}
\label{app:resampling}

Fix numbers \(x_1,\ldots,x_m\), let
\(\bar x=m^{-1}\sum_i x_i\), and draw \(I_1,\ldots,I_B\) independently and
uniformly from \([m]\). Define counts
\(C_i=\sum_{b=1}^B\I\{I_b=i\}\) and
\[
 \bar X^*=\frac1B\sum_{b=1}^B x_{I_b}
 =\frac1B\sum_{i=1}^m C_i x_i.
\]
Conditionally on the corpus,
\[
 (C_1,\ldots,C_m)\sim\operatorname{Multinomial}
 \left(B;\frac1m,\ldots,\frac1m\right),
 \qquad
 \E[\bar X^*\mid x_{1:m}]=\bar x.
\]
For \(B=m\), each \(C_i\sim\operatorname{Binomial}(m,1/m)\) and converges in
law to \(\operatorname{Poisson}(1)\) as \(m\to\infty\). Thus the expected
fraction of the corpus omitted by one bootstrap resample approaches
\(e^{-1}\). If
\(B=Em\) and \(E\) grows, then \(C_i/E\) concentrates near one: repeated
resampling reconstructs the original empirical weights.

Now let \(X_1,\ldots,X_m\) be square-integrable and possibly dependent, and set
\(Y_b=X_{I_b}\). Two resampled observations are independent conditional on
the corpus, but unconditionally
\begin{equation}
 \operatorname{Cov}(Y_1,Y_2)
 =\operatorname{Var}(\bar X)
 =\frac1{m^2}\sum_{r,s=1}^m\operatorname{Cov}(X_r,X_s).
 \label{eq:resampling-covariance}
\end{equation}
The law of total variance gives
\begin{equation}
 \operatorname{Var}(\bar X^*)
 =\operatorname{Var}(\bar X)
 +\frac1B\E\!\left[\frac1m\sum_{i=1}^m(X_i-\bar X)^2\right].
 \label{eq:resampling-variance}
\end{equation}
The second term is nonnegative. Averaging out the resampling randomness---the
Rao--Blackwell step---returns \(\bar X\). Resampling cannot erase the first
term, which contains the original dependence.

Small or structured subsamples can still change a statistical procedure. For
example, taking well-separated positions from a time series can lower pairwise
dependence, but it also uses fewer positions. Whether the tradeoff helps
depends on the estimator and the dependence structure; it is not a consequence
of replacement sampling itself.

\section{Proof of the independent-points comparison}
\label{app:mixing-proof}

\begin{proposition}[Independent-points comparison]
Under the assumptions preceding \cref{eq:independent-points},
\[
 \left\|\mathcal L(X_{t_1},\ldots,X_{t_q})-P^{\otimes q}\right\|_{\TV}
 \leq(q-1)\beta(a).
\]
\end{proposition}

\begin{proof}
Let \(Z_j=X_{t_j}\),
\(Q_j=\mathcal L(Z_1,\ldots,Z_j)\), and
\(d_j=\|Q_j-P^{\otimes j}\|_{\TV}\). Restricting the distributions in
\cref{eq:beta} to \((Z_1,\ldots,Z_j)\) gives
\[
 \|Q_j-Q_{j-1}\otimes P\|_{\TV}\leq\beta(a).
\]
Taking the product of both measures with the same probability law preserves
total variation distance. Since \(Q_1=P\), the triangle inequality gives
\[
 d_1=0,
 \qquad
 d_j\leq\beta(a)+d_{j-1}.
\]
Iteration yields \(d_q\leq(q-1)\beta(a)\).
\end{proof}

\clearpage
\begin{thebibliography}{99}

\bibitem[Bradley(2005)]{bradley2005}
Richard C. Bradley.
\newblock Basic properties of strong mixing conditions. A survey and some open
questions.
\newblock Probability Surveys, 2:107--144, 2005.
\newblock \url{https://doi.org/10.1214/154957805100000104}.

\bibitem[Ding et al.(2024)]{ding2024}
Hantian Ding, Zijian Wang, Giovanni Paolini, Varun Kumar, Anoop Deoras,
Dan Roth, and Stefano Soatto.
\newblock Fewer truncations improve language modeling.
\newblock 2024. \url{https://arxiv.org/abs/2404.10830}.

\bibitem[Dubey et al.(2024)]{dubey2024}
Abhimanyu Dubey, Abhinav Jauhri, Abhinav Pandey, and others.
\newblock The Llama 3 herd of models.
\newblock 2024. \url{https://arxiv.org/abs/2407.21783}.

\bibitem[Hairer(2002)]{hairer2002}
Martin Hairer.
\newblock Exponential mixing for a stochastic PDE driven by degenerate noise.
\newblock Nonlinearity, 15(2):271--279, 2002.
\newblock \url{https://arxiv.org/abs/math-ph/0103039}.

\bibitem[Hoeffding(1963)]{hoeffding1963}
Wassily Hoeffding.
\newblock Probability inequalities for sums of bounded random variables.
\newblock Journal of the American Statistical Association, 58(301):13--30,
1963. \url{https://doi.org/10.1080/01621459.1963.10500830}.

\bibitem[Lattimore and Szepesv\'ari(2020)]{lattimore2020}
Tor Lattimore and Csaba Szepesv\'ari.
\newblock \emph{Bandit Algorithms}.
\newblock Cambridge University Press, 2020.
\newblock \url{https://tor-lattimore.com/downloads/book/book.pdf}.

\bibitem[Penedo et al.(2024)]{penedo2024}
Guilherme Penedo, Hynek Kydl\'i\v{c}ek, Loubna Ben Allal, and others.
\newblock The FineWeb datasets: Decanting the web for the finest text data at
scale.
\newblock 2024. \url{https://arxiv.org/abs/2406.17557}.

\bibitem[Soldaini et al.(2024)]{soldaini2024}
Luca Soldaini, Rodney Kinney, Akshita Bhagia, and others.
\newblock Dolma: An open corpus of three trillion tokens for language model
pretraining research.
\newblock 2024. \url{https://arxiv.org/abs/2402.00159}.

\bibitem[Yu(1994)]{yu1994}
Bin Yu.
\newblock Rates of convergence for empirical processes of stationary mixing
sequences.
\newblock Annals of Probability, 22(1):94--116, 1994.
\newblock \url{https://doi.org/10.1214/aop/1176988849}.

\end{thebibliography}

\end{document}
